% GENERATED FILE. Edit canonical lessons, then run npm run generate:books.
% canonical-source-hash: fnv1a64:3be50bca3d66a70e
% canonical-lessons: AR-C295-hikaya, AR-C295-rawa, AR-C295-qassa, AR-C295-sarada, AR-C295-dhata-yawm

\chapter{Telling a story}
\label{ch:ar-telling-a-story}
\hlchaptermodality{295}

\begin{chapteropening}
\textbf{By the end of this chapter:} \emph{I can say a tale, to narrate, to tell a story, to recount and one day.}

Complete the last lesson of chapter 295: I can say a tale, to narrate, to tell a story, to recount and one day.
\end{chapteropening}

\section[\texorpdfstring{\ar{حكاية} (ḥikāya)}{حكاية (ḥikāya)}]{\ar{حكاية} (ḥikāya) — a tale, a story told aloud}

\label{lesson:AR-C295-hikaya}



\begin{quote}
Before the new one: say the Arabic for at first, then the Arabic for at that moment.
\end{quote}

\subsection*{You'll want to know: \ar{حكاية}}
\textbf{\ar{حكاية}} — \emph{ḥikāya} — \textquotedblleft{}a tale, a story told aloud\textquotedblright{}.

\begin{quote}
\ar{سمع} \ar{الأطفال} \ar{حكاية} \ar{قبل} \ar{النوم}.
\end{quote}

\begin{quote}
\emph{samiʿa al-aṭfāl ḥikāya qabla an-nawm.}
\end{quote}

\begin{quote}
The children heard a story before bed.
\end{quote}

Plural \ar{حكايات}. A \ar{حكاية} is told; a \ar{قصة} can be told or read.

\begin{cousinweb}
\textbf{\ar{حكاية}} (\emph{ḥikāya}) is built on the native Arabic root \emph{ḥ-k-y} (Lane), where \emph{ḥakā} means both \textquotedblleft{}to relate\textquotedblright{} and \textquotedblleft{}to imitate\textquotedblright{}: a tale repeats what someone said or did.
\end{cousinweb}

\begin{grammarlens}[title={What you've built}]
One more word for telling what happened.
\end{grammarlens}

\subsection*{Your turn}
\begin{itemize}
\raggedright
  \item \emph{Say it:} \emph{ḥikāya}
  \item \emph{Say it:} \emph{ḥikāya}, once more
  \item \emph{Say it:} \emph{samiʿa al-aṭfāl ḥikāya qabla an-nawm.} — The children heard a story before bed.
  \item \emph{Recall:} say the Arabic for at first, then the Arabic for at that moment, then say \emph{ḥikāya} again
\end{itemize}

\begin{tcolorbox}[breakable,colback=teal!4,colframe=teal!35!black,title={Before you move on}]
What does \ar{حكاية} mean? (\textquotedblleft{}A tale, a story told aloud\textquotedblright{}.) Say it once more.
\end{tcolorbox}
\section[\texorpdfstring{\ar{روى} (rawā)}{روى (rawā)}]{\ar{روى} (rawā) — to narrate, to relate}

\label{lesson:AR-C295-rawa}



\begin{quote}
Before the new one: say the Arabic for at that moment, then the Arabic for a tale.
\end{quote}

\subsection*{You'll want to know: \ar{روى}}
\textbf{\ar{روى}} — \emph{rawā} — \textquotedblleft{}to narrate, to relate\textquotedblright{}.

\begin{quote}
\ar{روى} \ar{جدي} \ar{قصة} \ar{عن} \ar{القرية}.
\end{quote}

\begin{quote}
\emph{rawā jaddī qiṣṣa ʿan al-qarya.}
\end{quote}

\begin{quote}
My grandfather told a story about the village.
\end{quote}

Present: \ar{يروي}, \emph{yarwī}.

The narrator is \ar{الراوي}; the taught \ar{رواية}, a novel, is the same root.

\begin{cousinweb}
\textbf{\ar{روى}} (\emph{rawā}) is built on the native Arabic root \emph{r-w-y} (Lane), which gives both \textquotedblleft{}to give water to drink\textquotedblright{} and \textquotedblleft{}to transmit a report or a poem\textquotedblright{}. \emph{Riwāya}, a novel, and \emph{rāwī}, a narrator, share it.
\end{cousinweb}

\begin{grammarlens}[title={What you've built}]
One more word for telling what happened.
\end{grammarlens}

\subsection*{Your turn}
\begin{itemize}
\raggedright
  \item \emph{Say it:} \emph{rawā}
  \item \emph{Say it:} \emph{rawā}, once more
  \item \emph{Say it:} \emph{rawā jaddī qiṣṣa ʿan al-qarya.} — My grandfather told a story about the village.
  \item \emph{Recall:} say the Arabic for at that moment, then the Arabic for a tale, then say \emph{rawā} again
\end{itemize}

\begin{tcolorbox}[breakable,colback=teal!4,colframe=teal!35!black,title={Before you move on}]
What does \ar{روى} mean? (\textquotedblleft{}To narrate, to relate\textquotedblright{}.) Say it once more.
\end{tcolorbox}
\section[\texorpdfstring{\ar{قص} (qaṣṣa)}{قص (qaṣṣa)}]{\ar{قص} (qaṣṣa) — to tell a story; to cut}

\label{lesson:AR-C295-qassa}



\begin{quote}
Before the new one: say the Arabic for a tale, then the Arabic for to narrate.
\end{quote}

\subsection*{You'll want to know: \ar{قص}}
\textbf{\ar{قص}} — \emph{qaṣṣa} — \textquotedblleft{}to tell a story; to cut\textquotedblright{}.

\begin{quote}
\ar{قصت} \ar{المعلمة} \ar{على} \ar{التلاميذ} \ar{قصة} \ar{قصيرة}.
\end{quote}

\begin{quote}
\emph{qaṣṣat al-muʿallima ʿalā at-talāmīdh qiṣṣa qaṣīra.}
\end{quote}

\begin{quote}
The teacher told the pupils a short story.
\end{quote}

Present: \ar{يقص}, \emph{yaquṣṣ}.

It takes \ar{على} for the listener: \ar{قص} \ar{على} \ar{الأطفال}. With scissors it means to cut.

\begin{cousinweb}
\textbf{\ar{قص}} (\emph{qaṣṣa}) is built on the native Arabic root \emph{q-ṣ-ṣ} (Lane), \textquotedblleft{}to follow a track step by step\textquotedblright{} and \textquotedblleft{}to clip\textquotedblright{}. To tell a story is to follow events in order; \emph{qiṣṣa}, a story, is its noun.
\end{cousinweb}

\begin{grammarlens}[title={What you've built}]
One more word for telling what happened.
\end{grammarlens}

\subsection*{Your turn}
\begin{itemize}
\raggedright
  \item \emph{Say it:} \emph{qaṣṣa}
  \item \emph{Say it:} \emph{qaṣṣa}, once more
  \item \emph{Say it:} \emph{qaṣṣat al-muʿallima ʿalā at-talāmīdh qiṣṣa qaṣīra.} — The teacher told the pupils a short story.
  \item \emph{Recall:} say the Arabic for a tale, then the Arabic for to narrate, then say \emph{qaṣṣa} again
\end{itemize}

\begin{tcolorbox}[breakable,colback=teal!4,colframe=teal!35!black,title={Before you move on}]
What does \ar{قص} mean? (\textquotedblleft{}To tell a story; to cut\textquotedblright{}.) Say it once more.
\end{tcolorbox}
\section[\texorpdfstring{\ar{سرد} (sarada)}{سرد (sarada)}]{\ar{سرد} (sarada) — to recount, to narrate in order}

\label{lesson:AR-C295-sarada}



\begin{quote}
Before the new one: say the Arabic for to narrate, then the Arabic for to tell a story.
\end{quote}

\subsection*{You'll want to know: \ar{سرد}}
\textbf{\ar{سرد}} — \emph{sarada} — \textquotedblleft{}to recount, to narrate in order\textquotedblright{}.

\begin{quote}
\ar{سرد} \ar{المسافر} \ar{تفاصيل} \ar{رحلته}.
\end{quote}

\begin{quote}
\emph{sarada al-musāfir tafāṣīl riḥlatihi.}
\end{quote}

\begin{quote}
The traveller recounted the details of his journey.
\end{quote}

Present: \ar{يسرد}, \emph{yasrud}.

The noun \ar{سرد} is narration, as in a novel.

\begin{cousinweb}
\textbf{\ar{سرد}} (\emph{sarada}) is built on the native Arabic root \emph{s-r-d} (Lane), first \textquotedblleft{}to stitch\textquotedblright{} and \textquotedblleft{}to link rings of mail one into the next\textquotedblright{}, then \textquotedblleft{}to carry on a story without a break\textquotedblright{}.
\end{cousinweb}

\begin{grammarlens}[title={What you've built}]
One more word for telling what happened.
\end{grammarlens}

\subsection*{Your turn}
\begin{itemize}
\raggedright
  \item \emph{Say it:} \emph{sarada}
  \item \emph{Say it:} \emph{sarada}, once more
  \item \emph{Say it:} \emph{sarada al-musāfir tafāṣīl riḥlatihi.} — The traveller recounted the details of his journey.
  \item \emph{Recall:} say the Arabic for to narrate, then the Arabic for to tell a story, then say \emph{sarada} again
\end{itemize}

\begin{tcolorbox}[breakable,colback=teal!4,colframe=teal!35!black,title={Before you move on}]
What does \ar{سرد} mean? (\textquotedblleft{}To recount, to narrate in order\textquotedblright{}.) Say it once more.
\end{tcolorbox}
\section[\texorpdfstring{\ar{ذات} \ar{يوم} (dhāta yawm)}{ذات يوم (dhāta yawm)}]{\ar{ذات} \ar{يوم} (dhāta yawm) — one day, once upon a time}

\label{lesson:AR-C295-dhata-yawm}



\begin{quote}
Before the new one: say the Arabic for to tell a story, then the Arabic for to recount.
\end{quote}

\subsection*{You'll want to know: \ar{ذات} \ar{يوم}}
\textbf{\ar{ذات} \ar{يوم}} — \emph{dhāta yawm} — \textquotedblleft{}one day, once upon a time\textquotedblright{}.

\begin{quote}
\ar{ذات} \ar{يوم} \ar{خرج} \ar{الصياد} \ar{إلى} \ar{البحر}.
\end{quote}

\begin{quote}
\emph{dhāta yawm kharaja aṣ-ṣayyād ilā al-baḥr.}
\end{quote}

\begin{quote}
One day the fisherman went out to sea.
\end{quote}

It opens a story, like the taught \ar{ذات} \ar{مرة}, once.

\begin{cousinweb}
\textbf{\ar{ذات} \ar{يوم}} (\emph{dhāta yawm}) puts \emph{yawm}, \textquotedblleft{}a day\textquotedblright{}, which is usually traced to an inherited Semitic word, after \emph{dhāta}, the feminine of \emph{dhū}, \textquotedblleft{}possessor of\textquotedblright{}, as in the taught \ar{ذات} \ar{مرة}.
\end{cousinweb}

\begin{grammarlens}[title={What you've built}]
One more word for telling what happened.
\end{grammarlens}

\subsection*{Your turn}
\begin{itemize}
\raggedright
  \item \emph{Say it:} \emph{dhāta yawm}
  \item \emph{Say it:} \emph{dhāta yawm}, once more
  \item \emph{Say it:} \emph{dhāta yawm kharaja aṣ-ṣayyād ilā al-baḥr.} — One day the fisherman went out to sea.
  \item \emph{Recall:} say the Arabic for to tell a story, then the Arabic for to recount, then say \emph{dhāta yawm} again
\end{itemize}

\begin{tcolorbox}[breakable,colback=teal!4,colframe=teal!35!black,title={Before you move on}]
What does \ar{ذات} \ar{يوم} mean? (\textquotedblleft{}One day, once upon a time\textquotedblright{}.) Say it once more.
\end{tcolorbox}
